% Vietnamese translation for OI Public Library
% Source-File: IOI中国国家候选队论文/2007/day2/9.周冬《生成树的计数及其应用》.ppt
\documentclass[11pt]{article}
\usepackage{oipl-vn}

\newcommand{\Oh}{\mathcal{O}}

\title{Bản trình chiếu: Đếm cây khung và ứng dụng}
\author{Zhou Dong}
\date{2007}

\begin{document}
\maketitle

\origtitle{Slide companion for counting spanning trees and applications}
\sourcepath{IOI China candidate team papers / 2007 / day 2 / Zhou Dong.ppt}

\section{Phạm vi bản dịch}

Tệp PowerPoint đi cùng bài viết đã được dịch từ PDF. Bản ghi chú này giữ lại
ý chính: từ bài toán đếm cây khung, đi qua định thức và định lý Matrix-Tree,
rồi quay lại các ứng dụng trong thi tin học.

\section{Bài toán}

Cho đồ thị vô hướng \(G\), cần tính số cây khung \(t(G)\). Với \(n\) nhỏ, có
thể dùng quy hoạch động theo tập với độ phức tạp mũ. Khi \(n\) lớn hơn, cần
một công cụ đại số.

\section{Định lý Matrix-Tree}

Lập ma trận Laplace của đồ thị: phần tử đường chéo là bậc của đỉnh, còn phần
tử ngoài đường chéo là âm số cạnh nối giữa hai đỉnh tương ứng. Xóa một hàng
và một cột bất kỳ khỏi ma trận này, định thức của ma trận con thu được chính
là số cây khung của \(G\).

\section{Tính định thức}

Trong cài đặt, ta tính định thức bằng khử Gauss hoặc biến thể nguyên nếu cần
tránh sai số. Với đồ thị có trọng số cạnh, cùng công thức cho tổng tích trọng
số của các cây khung, sau khi thay bậc bằng tổng trọng số kề và phần tử ngoài
đường chéo bằng âm tổng trọng số cạnh nối hai đỉnh.

\section{Ứng dụng}

Đếm cây khung xuất hiện trong bài toán mạng giao thông, xác suất chọn cây
ngẫu nhiên, đồ thị có nhiều cạnh song song, và các bài yêu cầu tính số cấu
trúc liên thông không chu trình. Matrix-Tree biến một bài đếm tổ hợp lớn
thành bài toán đại số kích thước \(n-1\).

\end{document}
